\section{Limitations and Future Work}
Several parts of the estimator still depend on calibrated runtime assumptions.
Runtime support and optimizer-step carryover are fudge terms that account for memory overhead outside the analytical model's current coverage.
Additionally, the slowdown score used for recommendations should become a measured timing model.

The next technical step is to connect timing and memory in the same artifact.
The memory estimator already ranks feasibility well enough for pre-launch
screening. A measured timing model would make slowdown estimates less
heuristic, especially for ZeRO-3, checkpointing, and tensor-parallel jobs where
communication and recomputation dominate user-visible cost.

\section{Related Work}
SimpleSFT follows prior work on memory-aware fine-tuning and distributed
training while targeting a narrower workflow: fast SFT strategy selection on
research LLMs before launch. LLMem estimates memory for distributed
fine-tuning and recommends methods for multi-GPU LLM training
\cite{kim2024llmem}, but it lacks modern attention backends and support for LoRA. Empirical work on distributed Transformer training
studies how model architecture and distributed strategy interact with memory
and performance \cite{lu2024distributed}. Broader resource-management work
compares analytical, CPU-side, and learning-based estimators for GPU memory and
utilization \cite{yousefzadeh2026gpu}. SimpleSFT uses the same general
motivation, then adds project-specific measurement artifacts and SFT-oriented
strategy export.

\begin{center}
\small
\begin{tabular}{@{}>{\raggedright\arraybackslash}p{0.31\linewidth}p{0.59\linewidth}@{}}
\toprule
Work & Relationship to SimpleSFT \\
\midrule
LLMem \cite{kim2024llmem} & Analytical GPU-memory estimation for LLM
fine-tuning and method recommendation across distributed training settings. It was calibrated using older models and different data regimes. \\
Transformer training study \cite{lu2024distributed} & Empirical view of
memory and performance behavior under distributed training choices. \\
Training-aware resource management \cite{yousefzadeh2026gpu} & Broader
resource-management framing for memory and utilization prediction. \\
\bottomrule
\end{tabular}
\normalsize
\end{center}

This work builds directly on prior work in training optimization. LoRA
reduces trainable state by injecting low-rank adapter matrices into frozen base
models \cite{hu2021lora}. ZeRO shards redundant training state across data
parallel ranks \cite{rajbhandari2020zero}. FlashAttention changes the memory
shape of attention through IO-aware tiling that avoids the full quadratic score
matrix \cite{dao2022flashattention}. These methods are common in current LLM
fine-tuning practice, including Qwen-family research workflows \cite{qwen25}.

\section{Conclusion}
SimpleSFT gives researchers a practical way to plan SFT jobs before reserving
scarce GPU time. The estimator connects inspected model structure with
distributed training choices, LoRA, checkpointing, FlashAttention, optimizer
state, and CUDA runtime support. The result is a phase-level estimate that can
guide resource selection and clarify the expected source of memory pressure.

The current saved measurement corpus shows a global mean absolute error of
1.332 GiB and a median relative error of 6.1 percent. Those numbers are strong
for pre-launch screening, especially when users keep several GiB of memory
headroom.
